\documentclass[10pt,a4paper]{article}

\usepackage[margin=1.05in]{geometry}
\usepackage{setspace}
\usepackage{amsmath,amssymb}
\usepackage{titlesec}
\usepackage[expansion=false]{microtype}
\usepackage{fancyhdr}

\setstretch{1.14}

\pagestyle{fancy}
\fancyhf{}
\fancyhead[L]{\small\itshape Elements of Position}
\fancyhead[R]{\small\itshape Overview}
\fancyfoot[C]{\small\thepage}
\renewcommand{\headrulewidth}{0pt}

\begin{document}

\begin{center}
\vspace*{2.0em}
{\large Elements of Position}\\[1.5em]
{\LARGE\bfseries Overview}\\[2.0em]
{\normalsize Justin Erholtz}\\[0.3em]
{\small 12 September 2026}
\end{center}

\vspace{1.8em}

\noindent
Where is \(1\)? Show me. Point to it.

\vspace{0.6em}

\noindent
You can't --- not because the mark is missing, but because one is an abstraction. We describe it.
One is a close: a unit that finishes itself. And in so doing holds a tension: at unit closure, the interior it encloses stays endlessly divisible.
That tension is where this project series lingers. 
It does not resolve the paradox so much as build outward from it, one part at a time.

Each part is proved before it is used, and trusted only once it has been run: built from something, checked against that thing twice, by two different routes, and shown before it is trusted --- never asserted from resemblance.
A later part does not explain an earlier one.
It stands on it.

Every step that can be proved is proved by basic geometry and algebra.
And in almost an odd way, at this stage none of this is claimed to be a new mathematical result, and the series says wherever it stands on established ground.

What is new is the figure those proofs assemble --- fixed in its internal ratios, indexed by the integers, the same figure at every scale.
That figure is treated with the care a load-bearing structure gets before anything is set on it: tested, named, and only trusted once it holds.

A working record sits alongside the numbered parts, revised as the object is measured further.
It is the walk that produces the pair --- two operators, a start on a diameter, a count of crossings --- without a finished turn inside the step.
What can be measured from that figure is documented and then proven from that walks trace.
What is only a count is reported from its tables, and named as a count, not smuggled in as an identity.
A length that appears only after those tables exist is named as a comparison.

This work tries to hold two things at once that are usually treated as incompatible.
The first is respect: for the disciplines it moves through, for the work already done in them, and for the thinkers reaching back into deep antiquity who asked versions of these questions with the tools they had.
Nothing here is written to outflank or dismiss that lineage.
The second is a refusal to be tidy about it.
Where mathematics, physics, architecture, sound, and philosophy touch --- and they touch more often, and more strangely, than any one of them is comfortable admitting alone --- this series lingers there rather than picking a side and moving on.

An idea is not owed patience for being old, and it is not owed patience for being recent, or for being dense.
What is asked here is what is asked of the figure itself: that a claim be built from something, checked, and shown, before it is named.
The series does not commit to a single open problem, a single field, or a fixed endpoint.
It commits to building one figure honestly, honoring the history that made the question askable, and staying in the overlaps rather than forcing them shut.

\end{document}
